\documentclass[11pt]{report}

%packages
\usepackage{geometry,
amsmath,
amssymb,
amsthm,
setspace,
mathrsfs,
tikz}

%This is to give color.
\usepackage{xcolor}

%This is to give space on the sides
%and between lines for comments
\geometry{margin=1.in}
%\doublespacing

\theoremstyle{plain}
\newtheorem{thm}{Theorem}
\newtheorem{lem}[thm]{Lemma}
\newtheorem{prop}[thm]{Proposition}
\newtheorem{cor}[thm]{Corollary}

%%%%
%% New Commands
%%%%
\newcommand{\modn}[3]{#1 \equiv #2 \: (\text{mod } #3)}
\newcommand{\Z}{\mathbb{Z}}
\newcommand{\R}{\mathbb{R}}
\newcommand{\Q}{\mathbb{Q}}
\newcommand{\N}{\mathbb{N}}
\newcommand{\bs}{$\backslash$}
\newcommand{\ps}{\mathscr{P}}

%%%%%%%
% Document Begins
%%%%%%%%
\begin{document}
\hfill Math 265

%%% Your Name goes HERE
\hfill \textbf{YOUR NAME HERE}

\hfill \today

\begin{center}
\Large{\textbf{Homework \# 13}} \\

\end{center}

\noindent\fbox{Problem List} \\ \S 12.5 \# 2, 6, 8, 9\\
\S 12.6 \# 2, 5, 6, 7, 8\\
\S 14.1 \# 2, 8, 10, 11, 14 \\
\S 14.2 \# 2, 4, 5, 7, 8, 11, 12\\

\begin{enumerate}
%problem
\item[\textbf{12.5.2}] In Exercise 9 of Section 12.2 you proved that $f : \R -\{2\} \to \R - \{5\} $ defined by $f (x) = \frac{5x+1}{x-2}$ is bijective. Now find its inverse.

\textbf{Answer:} \\

\vspace{1in}

%problem
\item[\textbf{12.5.6}] The function $f : \Z \times \Z \to \Z \times \Z$  defined by the formula $f (m,n) = (5m+4n,4m+3n)$ is bijective. Find its inverse.

\textbf{Answer:} \\

\vspace{1in}

%problem
\item[\textbf{12.5.8}] Is the function $\theta : \mathcal{P}(\Z) \to \mathcal{P}(\Z)$ defined as $\theta(X) = \overline{X}$ bijective? If so, find $\theta^{-1}.$

\textbf{Answer:} \\

\vspace{1in}

%problem
\item[\textbf{12.5.9}] Consider the function $f : \R \times \N \to \N \times \R$ defined as $f (x, y) = (y,3x y).$ Check that this
is bijective; find its inverse.

\textbf{Answer:} \\

\vspace{1in}

%problem
\item[\textbf{12.6.2}] Consider the function $f : \{1,2,3,4,5,6,7\} \to \{0,1,2,3,4,5,6,7,8,9\}$ given as $f =\{(1,3), (2,8), (3,3), (4,1), (5,2), (6,4), (7,6)\}.$  Find 
	\begin{enumerate}
	\item[a.] $f(\{1,2,3\})$
	
	\textbf{Answer:} \\
	
	\item[b.] $ f (\{4,5,6,7\}$
	
	\textbf{Answer:} \\
	
	\item[c.] $f(\emptyset)$
	
	\textbf{Answer:} \\
	
	\item[d.] $f^{-1}(\{0,5,9\})$
	
	\textbf{Answer:} \\
	
	\item[e.] $f^{-1}(\{0,3,5,9\})$
	
	\textbf{Answer:} \\
	\end{enumerate}
	
\vspace{1in}

%problem
\item[\textbf{12.6.5}] Consider a function $f : A \to B$ and a subset $X \subseteq A.$ We observed in Example 12.14 that $f^{-1}(f (X)) \not = X$ in general. However $X \subseteq f^{-1}(f (X))$ is always true. Prove this.

\textbf{Answer:} \\

\vspace{1in}

%problem
\item[\textbf{12.6.6}] Given a function $f : A \to B$ and a subset $Y \subseteq B$, is $f (f(Y)) = Y$ always true? Prove or give a counterexample.

\textbf{Answer:} \\

\vspace{1in}

%problem
\item[\textbf{12.6.7}] Given a function $f : A \to B$ and subsets $W, X \subseteq A,$ prove $f (W \cap X) \subseteq f (W)\cap f (X).$

\textbf{Answer:} \\

\vspace{1in}

%problem
\item[\textbf{12.6.8}] Given a function $f : A \to B$ and subsets $W, X \subseteq A,$ then $f (W \cap  X) = f (W) \cap f (X)$ is false in general. Produce a counterexample.

\textbf{Answer:} \\

\vspace{1in}

\textbf{Directions for \S 14.1 problems:} Show that the two given sets have equal cardinality by describing a bijection
from one to the other. Describe your bijection with a formula (not as a table).\\
%problem
\item[\textbf{14.1.2}] $\R$ and $(\sqrt{2},\infty)$

\textbf{Answer:} \\

\vspace{1in}

%problem
\item[\textbf{14.1.8}] $\Z$ and $S=\{x \in \R : \sin(x)=1\}$

\textbf{Answer:} \\

\vspace{1in}

%problem
\item[\textbf{14.1.10}] $\{0,1\} \times \N$ and $\Z$

\textbf{Answer:} \\

\vspace{1in}

%problem
\item[\textbf{14.1.11}] $[0,1]$ and $(0,1)$\\

\textbf{Comment:} This is very tricky. I gave a hint at the end of the homework. Also, we will learn an easier way to handle this in the future.

\textbf{Answer:} \\

\vspace{1in}

%problem
\item[\textbf{14.1.14}]  $\N \times \N$ and $\{(n,m) \in \N \times \N : n \leq m\}$

Hint: I think it's easier to find a bijection from $\N \times \N$  to $\{(n,m) \in \N \times \N : n \leq m\}$. One way to help ensure the injectivity is to keep one coordinate the same.

\textbf{Answer:} \\

\vspace{1in}
\textbf{Directions:} For the problems in section 14.2, you may use any theorem \emph{from this section.} You must reference it the theorem explicitly as in ``Applying Theorem 14.5 that says....., we can conclude...."\\

%problem
\item[\textbf{14.2.2}] Prove that the set $A =\{(m,n) \in \N \times \N : m \leq  n\}$ is countably infinite.

\textbf{Answer:} \\

\vspace{1in}

%problem
\item[\textbf{14.2.4}] Prove that the set of all irrational numbers is uncountable. (Suggestion: Consider proof by contradiction using Theorems 14.4 and 14.6.)

\textbf{Answer:} \\

\vspace{1in}

%problem
\item[\textbf{14.2.5}] Prove or disprove: There exists a countably infinite subset of the set of irrational numbers.

\textbf{Answer:} \\

\vspace{1in}

%problem
\item[\textbf{14.2.7}] Prove or disprove: The set $\Q^{100}$ is countably infinite.

\textbf{Answer:} \\

\vspace{1in}

%problem
\item[\textbf{14.2.8}] Prove or disprove: The set $\Z\times \Q$ is countably infinite.

\textbf{Answer:} \\

\vspace{1in}

%problem
\item[\textbf{14.2.11}] Describe a partition of $\N$ that divides $\N$ into eight countably infinite subsets.

\textbf{Answer:} \\

\vspace{1in}

%problem
\item[\textbf{14.2.12}] Describe a partition of $\N$ that divides $\N$ into $\aleph_0$ countably infinite subsets.

\textbf{Answer:} \\

\vspace{1in}


\textbf{Hint to 14.1.11} One way to think about finding a bijection $f: (0,1) \to [0,1]$ is that we just need to fill two ``holes" at $x=0$ and at $x=1.$ But if I define $f(1/2)=1$, how do we replace it? \\

Consider the infinite sets $S=\{1/2, 1/3, 1/4, 1/5, \cdots \}$ and $T=\{1/3, 1/4, 1/5, \cdots \}.$ How could you show that $|S|=|T|$? How can you use this idea to fill the hole at $x=1/2$?
\end{enumerate}
\end{document}